\section{The Thinking of Sound and Well-Born Persons}

In his self-defense, Evola denied he was a Fascist and insisted that he thought the same as well-bred men with sound minds thought prior to the French Revolution. This recording by \textbf{Solange Hertz} remarkably illustrates how such men thought and viewed the world. She identifies the ``Three Plagues'' of the modern mind, which will be familiar to any who studies Guenon or Evola.

\begin{enumerate}
\item Scientism
\item Liberalism
\item Universalism
\end{enumerate}

The Three Plagues of the Great Apostasy!\footnote{\url{https://gornahoor.net/library/ThreePlagues.mp3}}

Whatever your opinion about the Roman Church, it is instructive to hear a Traditional voice, speaking as though the past 400 years never happened.

Note: Originally appears at \textit{In the Spirit of Chartres}\footnote{\url{https://isoc.ws/}}.


\flrightit{Posted on 2010-09-17 by Cologero}

\begin{center}* * *\end{center}

\begin{footnotesize}
\begin{sffamily}
\texttt{man on 2010-09-17 at 18:26 said:}

Breeding is for dogs not thinkers.

\hfill

\texttt{Cologero on 2010-09-17 at 19:22 said:}

Unfortunately, you know nothing of aristocracy nor of Evola. This is the exact quote:

\begin{quote}\footnotesize
Perciò, nei miei riguardi, di apologia di ``idee proprie al fascismo'' non è affatto il caso di parlare. I miei principi sono solo quelli che prima della Rivoluzione francese ogni persona ben nata considerava sani e normali.
\end{quote}


\hfill

\texttt{Andrew on 2010-09-19 at 05:30 said:}

I find myself delighted whenever I read or hear Traditionalist Catholic material, but the one thing that always bothers me is the hurdle of Perennialism / Tradition. I would be quite in agreement and seek the association of such persons if it were not for the exclusivist attitude in Christianity.

\hfill

\texttt{Cologero on 2010-09-19 at 11:06 said:}

Keep in mind that the perspective of a Guenon or Evola --- who regarded themselves as beyond caste or specific tradition --- is not for everyone. Most people will always follow their own tradition, and rightfully so. An authentic Tradition is complete in itself and does not need to even know about other Traditions. Nevertheless, early Catholic theology absorbed other traditions --- in an authentic, not a syncretic way. Joseph de Maistre regarded the Catholic Faith as the primordial tradition ... that is an idea that needs developing within the Church.

An interesting point she made was that so-called primitive peoples are actually degenerations of a higher civilisation, not an earlier stage of it. This is straight from Guenon and Evola. If someone can identify a Catholic source for that idea, please let us know.

As for exclusivity, we can see its positive aspects as loyalty to one's own tradition and the setting of boundaries against outside and perverse influences. The danger of perennialism is that the attempt to be ``everything'' will result in being nothing in particular. Both Guenon (Islam) and Evola (paganism) had to settle the personal question in a particular way.

\hfill

\texttt{man on 2010-09-19 at 15:51 said:}

Cologero, ``unfortunately'' my knowledge and understanding Evola is not something you have enough data to judge. In any case I do not share all his views. Tell me, do you consider yourself a well bred aristocrat?

\hfill

\texttt{Cologero on 2010-09-19 at 16:37 said:}

I took your comment into consideration and retranslated the title to be more consistent with Evola's actual words and intent.

\hfill

\texttt{man on 2010-09-19 at 17:42 said:}

I will interpret your avoiding the question as a no. Good job on the blog by the way.

\hfill

\texttt{Mark Eppihimer on 2010-09-22 at 22:52 said:}

I really enjoyed the recording, and find myself wondering where I may be today if I had heard such fascinating ideas on Traditional Catholicism when I struggled with the faith so many years ago.

\hfill

\texttt{GFJF on 2010-09-29 at 23:37 said:}

Man: In reference to men, `well bred' means well brought up, and has nothing to do with eugenic breeding of biological traits.

Charles: I don't speak Italian but my guess is that well bred renders Evola's meaning better than well born. He probably meant his statement to include faithful serfs and freemen as well as lords. It's worth noting that Plato also stresses good breeding, and is at pains to show that good birth doesn't guarantee virtue.

\hfill

\texttt{Bryan M. M. Reynolds on 2017-11-02 at 22:15 said:}

primitive peoples are actually degenerations of a higher civilisation, not an earlier stage of it

If someone can identify a Catholic source for that idea, please let us know.

Wolfgang Smith mentions Patristic descriptions of Paradise before the Fall. For example: ``As St. John Chrysostom observes: `Man lived on earth like an Angel; he was in the body, but had no bodily needs.' One sees that by contemporary standards the Age of Paradise is still miraculous; only after the Fall --- and by virtue of the Fall --- do the contours of our world begin to come into view, for indeed, `By man came death.' (1 Cor. 1 5.21)''

... and:---


St. Macarius the Great informs us that the bodily expulsion from Paradise had its counterpart in the soul: ``That Paradise was closed,'' he writes, ``and that a Cherubim was commanded to prevent men from entering it by a flaming sword: of this we believe that in visible fashion it was just as it is written, and at the same time we find that this occurs mystically in every soul.'' It was then that the ``carnal man,'' the psychikos anthropos ``who receiveth not the things of the Spirit of God: for they are foolishness unto him,'' and who knows them not ``because they are spiritually discerned'' (1 Cor. 2. 14) --- it was then that man as we know him came into being. Not instandy, in fact, but gradually; for as we also learn from Genesis, Adam and Eve remained for some time ``close'' to Paradise: close enough to see it from afar. It can further be said, moreover, that mankind has been engaged since primordial times in an ongoing fall from Paradise: every betrayal, large or small --- everything that theology knows as ``sin'' --- constitutes a link, as it were, in this (decidedly non-Darwinist!) ``chain of descent.''

Smith recommends Seraphim Rose's \_Genesis, Creation and Early Man\_ (Platina, CA, 2000) as an indispensable roundup.


\hfill

\end{sffamily}
\end{footnotesize}

